% Capítulo 10 – gerado a partir do relatório; edite livremente
\chapter{CONCLUSÕES}

O planejamento apresentado permite fabricar a engrenagem do PROTEC a partir de uma barra de SAE 8620 de Ø160 $\times$ 58 mm em onze operações, com cinco processos de usinagem: torneamento, furação, fresamento, brochamento e retificação. A escolha de gerar o furo e a face do cubo na mesma fixação, e de retificar o furo no final com a peça centrada pelo diâmetro primitivo, garante a concentricidade entre o dentado e o furo e a perpendicularidade das faces de encosto, que são as características mais importantes para o funcionamento da engrenagem.


A verificação das potências mostrou que, para uma peça de Ø160 mm, o torno Centur 30D trabalha limitado pelo torque, e os parâmetros de desbaste precisaram ser reduzidos em relação aos valores nominais dos catálogos. As demais máquinas trabalham com folga. O fresamento dos dentes com fresa de disco é adequado a um lote pequeno, mas responde pela maior parte do tempo de usinagem; em uma produção seriada, a geração com fresa caracol reduziria esse tempo de forma expressiva.

